% Hand-written narrative; numbers are macros.
Table~\ref{tab:ablation} changes one component at a time on the same test windows (seeds 0--2).
The per-user baseline (C1) carries the method: replacing it with a population baseline drops the
detection rate from \AblFullDR{} to \AblNoCOneDR{} and the precision from \AblFullP{} to
\AblNoCOneP{}, the level of Isolation Forest. Removing the resource family lowers the detection rate
to \AblNoResourceDR{}, and removing the volume family lowers the precision to \AblNoVolumeP{}.
Repeat-alert suppression (C4), switched off by tuning, lowers precision (\AblCFourP).

Several variants do \emph{better} on the test windows than the tuned model: Fisher fusion
(DR \AblFisherDR, P \AblFisherP), no calibration (DR \AblNoCTwoDR, P \AblNoCTwoP) and dropping the
identity family (DR \AblNoIdentityDR, P \AblNoIdentityP). Validation had preferred the tuned
settings (validation DR@10 \ValDROurs{} for the selected configuration). We do not switch to any of
these variants, because choosing among them now would select on the test windows; we report the gap
as evidence that the short development period, whose validation windows hold fewer insiders than
the test windows, is a weak guide to the later months.

\textbf{Fusion method.} Validation selected Bonferroni min-$p$ fusion over Fisher and Stouffer
combinations. Min-$p$ is conservative: it is driven by whichever behaviour family produced the
most extreme $p$-value, which in practice means the family with the largest $\ell_2$ deviation.
Fisher's method, which combines evidence across all four families, performs better on the test
windows, suggesting that true threats in this dataset tend to deviate in multiple families
simultaneously. This is consistent with S2 insiders, whose job-site browsing (resource family)
and unusual USB activity (volume family) often co-occur on the same day.

\textbf{Calibration.} Removing pooled conformal calibration (C2) and replacing it with a
rank-based score improves the test detection rate. This is surprising: calibration was chosen on
validation, but the score distributions differ between the development period and the later test
months, so calibration thresholds that were correctly set on development data are slightly
miscalibrated on later data. The paper reports this openly as evidence of the distribution shift
noted throughout the evaluation.

\begin{table}[t]
\centering
\caption{Ablation at 10 alerts/day on the test windows (DR with 95\% CI over insiders; P@10).}
\label{tab:ablation}
\resizebox{\columnwidth}{!}{%
\begin{tabular}{lcc}
\toprule
Variant & DR@10 [95\% CI] & P@10 \\
\midrule
Full (tuned) & 0.43 [0.28, 0.60] & 0.028 \\
-C1 population baseline & 0.28 [0.14, 0.44] & 0.007 \\
-C2 no calibration & 0.63 [0.47, 0.79] & 0.021 \\
-C3 no fusion (pooled) & 0.56 [0.39, 0.72] & 0.019 \\
+C4 suppress 3 d & 0.43 [0.27, 0.58] & 0.018 \\
-identity family & 0.55 [0.39, 0.71] & 0.034 \\
-resource family & 0.30 [0.17, 0.45] & 0.029 \\
-temporal family & 0.49 [0.34, 0.66] & 0.028 \\
-volume family & 0.39 [0.24, 0.57] & 0.009 \\
fusion: fisher & 0.61 [0.45, 0.76] & 0.025 \\
fusion: harmonic & 0.49 [0.33, 0.65] & 0.029 \\
fusion: stouffer & 0.17 [0.06, 0.30] & 0.006 \\
Isolation Forest (reference) & 0.29 [0.15, 0.45] & 0.007 \\

\bottomrule
\end{tabular}}
\end{table}
